% Section 9: Conclusion

\section{Conclusion}

This paper measured threat-weighted cooperative and non-cooperative coverage as separate quantities against CapEx, under ray-traced line of sight on OSM buildings, on a synthetic city, an airport, and one real corridor. Cooperative coverage saturates with RF-inclusive mixes at modest CapEx: above \(0.96\) in every city run, and \(0.93\) for \$150k on Avenida Paulista. Non-cooperative coverage stays incomplete: \(0.41\)--\(0.51\) across three city seeds and \(0.22\) on Paulista. The same stored plans cover the 20\,m airway far worse than the 45--65\,m airways. Matched-CapEx controls show that RF alone suffices for lean cooperative coverage and that mixed Radar/EO/Acoustic sets give the best non-cooperative coverage only near \$320k, with or without site fees.

The fronts are reachable trade-offs, not optimal designs. At 608 evaluations, uniform random search reaches the same city hypervolume as NSGA-II and shows the same asymmetry, so the finding concerns the coverage layers, not the search procedure. A BVLOS safety case should therefore state both layers, or a dual-layer front with explicit floors, rather than one fused coverage percentage.

Section~\ref{sec:limitations} lists the limits on these findings; the next steps are diffraction-aware propagation, additional OSM morphologies, full-pool modality ablations, field calibration of \(P_D\), and finer voxels for low airways. Code: \url{https://github.com/andrekuros/SCOPAS_Sensor_Placement}.
